\section{Lifting and Intermediate Representation}
\label{sec:lifting}

\subsection{A strict decoder}

The decoder is not a general-purpose disassembler. It recognizes instruction
lengths, prefixes, ModRM/SIB fields, and operand structure only for the supported
subset. Every instruction is decoded from the function entry, and branch targets
must land on instruction boundaries inside the selected range. Two instructions
may not overlap. A relative branch outside the range, an unsupported prefix, an
unknown opcode, a malformed ModRM, or a truncated instruction is a build error
with an offset and a reason.

This strictness is a security property: there is no path on which an unsupported
instruction is skipped, approximated, or interpreted with the wrong semantics.
The decoder is also intentionally narrow in its prefix handling. Segment
overrides that do not affect 64-bit addressing are accepted and ignored so that
compiler-emitted \code{cs nop} padding lifts. Address-size overrides, lock
prefixes, and unsupported repeated prefixes fail closed. The undefined 16-bit
\code{bswap} form with a \code{0x66} operand-size prefix is rejected because real
CPUs implement it inconsistently.

\subsection{Virtual register file}

The VM has 96 64-bit virtual registers. Table~\ref{tab:registers} lists the
mapping. Architectural GPRs occupy slots 0--15, flags occupy 16--19, constant
zero and one occupy 20--21, compiler scratch occupies 24--59, and lifter scratch
occupies 60--63. The sixteen XMM registers occupy slots 64--95, two 64-bit halves
per register.

\begin{table}[htbp]
\centering
\small
\caption{Virtual register map.}
\label{tab:registers}
\begin{tabular}{@{}lll@{}}
\toprule
Slots & Contents & Notes \\
\midrule
0--15 & \code{RAX}, \code{RCX}, \code{RDX}, \code{RBX}, \code{RSP}, \code{RBP}, \code{RSI}, \code{RDI}, \code{R8}--\code{R15} & architectural GPRs \\
16--19 & \code{ZF}, \code{CF}, \code{SF}, \code{OF} & eagerly computed flags \\
20--21 & zero, one & constants \\
24--59 & compiler scratch & temporaries \\
60--63 & lifter scratch & expression lowering \\
64--95 & \code{XMM0}--\code{XMM15} & two 64-bit halves each \\
\bottomrule
\end{tabular}
\end{table}

Registers are masked at rest with a per-invocation key. A read unmasks only the
operand consumed by the current instruction; a write stores the result masked.
The key is derived from the effective payload key and is present in the generated
code. This is an at-rest masking layer, not cryptographic isolation from an
attacker who executes or debugs the VM.

\subsection{Flag model}

Flags are computed eagerly after every supported flag-setting instruction. The
lifter emits explicit comparisons into virtual flag registers, so a subsequent
\code{jcc}, \code{cmovcc}, or \code{setcc} reads a virtual register rather than
reconstructing native flag state.

\begin{itemize}
  \item \code{inc} and \code{dec} preserve \code{CF}.
  \item Shifts and \code{imul} mark flags unknown. A dependent conditional branch
        is rejected until a supported flag setter overwrites the unknown flag.
  \item Flag knowledge resets at every basic-block boundary. Cross-block flag
        liveness that cannot be proven is a build error, not an assumption.
  \item Parity conditions are unsupported; the decoder rejects them.
\end{itemize}

This design trades code size for determinism. It also makes the virtual
instruction stream independent of the host's native flag encoding.

\subsection{Private memory model}

Virtual memory is a fixed 4 KiB byte-addressable stack. The initial stack pointer
is 4080, the stack grows down, and a fixed read-only data blob of at most 2 KiB is
installed below the stack. Effective addresses are computed from base, index,
scale, and displacement using 64-bit wrapping arithmetic. Every access checks
\code{address + width <= memory.len()}; an out-of-range access fails closed.

RIP-relative reads of read-only data are resolved at build time. The referenced
bytes are copied into the private data blob, and the lifted instruction reads
from that private copy. Writable RIP-relative targets and absolute host addresses
are rejected. The VM never dereferences a host pointer.

\subsection{IR and CFG construction}

The IR is a fixed-capacity, two-address representation. Each basic block has one
terminator: a jump, a conditional branch, a return, a trap, or a direct call.
Leaders are discovered from the function entry, branch targets, and conditional
fallthroughs. Blocks are stored in a bounded vector; the function limit is
16 KiB, the block limit is 256, and the IR operation limit is 8192.

The lifter performs three consistency checks before flattening:

\begin{enumerate}
  \item every decoded instruction belongs to exactly one block;
  \item every branch target is a leader inside the function;
  \item every block ends in exactly one supported terminator and does not fall
        through accidentally.
\end{enumerate}

The resulting CFG is acyclic or cyclic, but it is always explicit. There are no
computed jump targets and no jump tables. This is what makes flattening and
fixed-pass execution possible.

\subsection{Supported x86-64 extensions}

Beyond the core integer subset, the lifter supports:

\begin{itemize}
  \item \textbf{RIP-relative read-only data}: a \code{lea} or memory operand with
        a RIP-relative displacement is resolved to a private data offset if the
        target lies in a read-only section and the bytes fit the data budget.
  \item \textbf{SSE2 integer operations}: \code{movdqu}/\code{movdqa}, bitwise
        \code{pxor}/\code{pand}/\code{por}, \code{paddq}/\code{psubq},
        \code{paddd}/\code{psubd}, \code{pmuludq},
        \code{punpcklqdq}/\code{punpckhqdq}, \code{psrldq}/\code{pslldq}, and
        \code{movq}/\code{movd} forms.
  \item \textbf{VEX-encoded 128-bit AVX integer forms} for the same operations,
        including \code{vpsrldq}/\code{vpslldq} and \code{vmovq}. AVX-256 and
        EVEX are rejected.
  \item \textbf{Byte swap}: 32-bit and 64-bit \code{bswap}, in normal and
        constant-time modes.
  \item \textbf{Repeated strings}: \code{rep movsb}/\code{movsq} and
        \code{rep stosb}/\code{stosq} over private memory in normal mode.
        Constant-time repeated strings fail closed.
  \item \textbf{Direct calls}: \code{call rel32} to another selected function,
        with a bounded guest call stack, and optional automatic callee following
        for resolvable local targets.
\end{itemize}

Table~\ref{tab:x86subset} in Appendix~\ref{app:x86} lists the accepted forms by
category. The list is intentionally conservative: an instruction that is not
modeled is a build error even if a native CPU would execute it.

\subsection{Rejected forms and the fail-closed contract}

The following are rejected at build time:

\begin{itemize}
  \item external calls and unresolved direct calls;
  \item indirect calls and indirect jumps;
  \item writable RIP-relative globals and absolute memory;
  \item jump tables and computed targets;
  \item 256-bit AVX/AVX2, EVEX, x87, MMX, and floating point;
  \item syscalls, \code{rdtsc}, \code{cpuid}, privileged or I/O instructions;
  \item atomics and fences;
  \item segment overrides that affect addressing, address-size overrides, and
        lock prefixes;
  \item cross-block flag liveness that cannot be proven;
  \item functions larger than 16 KiB or programs exceeding fixed capacities.
\end{itemize}

A build error is not a silent downgrade. The protector reports the offset and
reason, and the output file is not produced. This is a deliberate usability
trade-off: the supported subset is smaller than a commercial protector's, but a
successful build does not contain a mistranslated instruction.
